% E1 results (RQ1; C1, C2, C6). Every number is a res-macro key defined in values.tex
% (export_latex.py) or extra_values.tex (local/extra_values.py). Facts that are not numbers (which
% pages lack a title or language after repair, the kinds of regressions, which task reached the
% output cap, the audit loops) were checked in runs/e1/scores/*.jsonl and runs/e1/artifacts
% (2026-11-14). Differences are <metric>-diff-local-qwen25vl-7b-local-qwen3-8b = QV minus Q3.
\subsection{RQ1: Capability profiles (C1, C2, C6)}\label{sec:rq1}
Table~\ref{tab:main} gives each model's scores. No model has all six: \Q{} reads no images (C4)
and \QV{} ran no E2 session, so no CCS is reported. Differences below are \QV{} minus \Q{} on the
same resamples; we call the two models different only when that interval excludes~0.

\textbf{Generation (C2, C6).} \res{functional-gen-local-qwen3-8b-count} of the
\res{functional-gen-local-qwen3-8b-n} pages of \Q{} and \res{functional-gen-local-qwen25vl-7b-count}
of the \res{functional-gen-local-qwen25vl-7b-n} pages of \QV{} passed their task's functional
checks. With the accessibility sentence, \Q{} reached $S_\text{gen}={}$\res{sgen-local-qwen3-8b}
(95\,\% CI \res{sgen-local-qwen3-8b-ci}) and \QV{} \res{sgen-local-qwen25vl-7b} (CI
\res{sgen-local-qwen25vl-7b-ci}); without it, \res{sgen-unguided-local-qwen3-8b} and
\res{sgen-unguided-local-qwen25vl-7b}. Neither difference is distinguishable from 0 (nudged
\res{sgen-diff-local-qwen25vl-7b-local-qwen3-8b}, CI \res{sgen-diff-local-qwen25vl-7b-local-qwen3-8b-ci};
unguided \res{sgen-unguided-diff-local-qwen25vl-7b-local-qwen3-8b}, CI
\res{sgen-unguided-diff-local-qwen25vl-7b-local-qwen3-8b-ci}). The models differ in how they respond
to the sentence. It raised \Q's score by \res{gap-local-qwen3-8b} (CI \res{gap-local-qwen3-8b-ci})
and cut its violating nodes per 100 elements from \res{vp100-unguided-local-qwen3-8b} to
\res{vp100-nudged-local-qwen3-8b}, but made no detectable difference for \QV{} (default gap
\res{gap-local-qwen25vl-7b}, CI \res{gap-local-qwen25vl-7b-ci}; difference of the gaps
\res{gap-diff-local-qwen25vl-7b-local-qwen3-8b}, CI \res{gap-diff-local-qwen25vl-7b-local-qwen3-8b-ci}).
H5, that most models gain from the sentence, holds for \res{h5-count} models and is therefore
\res{h5-result} (Holm-adjusted $p={}$\res{h5-padj}). A high pass rate did not mean a clean page:
a share of only \res{clean-nudged-local-qwen3-8b} of \Q's nudged pages and
\res{clean-nudged-local-qwen25vl-7b} of \QV's passed every applicable check (difference
\res{clean-nudged-diff-local-qwen25vl-7b-local-qwen3-8b}, CI
\res{clean-nudged-diff-local-qwen25vl-7b-local-qwen3-8b-ci}).
\res{capped-gen-local-qwen3-8b-count} of \Q's \res{capped-gen-local-qwen3-8b-n} answers, all for the
icon-links task, reasoned until the output cap and contained no page; they score~0. Asked whether
its page conformed to WCAG~2.2 AA, \Q{} said yes for a share of \res{overclaim-local-qwen3-8b} (CI
\res{overclaim-local-qwen3-8b-ci}) of its functional pages that failed at least one check, and \QV{}
for \res{overclaim-local-qwen25vl-7b} (CI \res{overclaim-local-qwen25vl-7b-ci}; difference
\res{overclaim-diff-local-qwen25vl-7b-local-qwen3-8b}, CI
\res{overclaim-diff-local-qwen25vl-7b-local-qwen3-8b-ci}). These pooled rates hide where the claims
came from: almost all were made for nudged pages. On failing nudged pages the rate was
\res{overclaim-nudged-local-qwen3-8b} for \Q{} (CI \res{overclaim-nudged-local-qwen3-8b-ci}) and
\res{overclaim-nudged-local-qwen25vl-7b} for \QV{} (CI \res{overclaim-nudged-local-qwen25vl-7b-ci}),
on failing unguided pages \res{overclaim-unguided-local-qwen3-8b} and
\res{overclaim-unguided-local-qwen25vl-7b}. Asking for accessibility thus made both models far more
likely to claim it, whether or not the page improved. The checks find only some failures, so
these rates are lower bounds.

\textbf{Repair (C2).} Repair was much harder than generation, and here the models differ clearly.
\Q{} reached $S_\text{rep}={}$\res{srep-local-qwen3-8b} (CI \res{srep-local-qwen3-8b-ci}) and
\QV{} \res{srep-local-qwen25vl-7b} (CI \res{srep-local-qwen25vl-7b-ci}; difference
\res{srep-diff-local-qwen25vl-7b-local-qwen3-8b}, CI \res{srep-diff-local-qwen25vl-7b-local-qwen3-8b-ci}).
Full repairs were rare: a share of \res{fullrepair-local-qwen3-8b} of \Q's pages and none of
\QV's (\res{fullrepair-local-qwen25vl-7b}).
The two models failed in opposite ways. \QV{} often changed nothing:
\res{unchanged-repair-local-qwen25vl-7b-count} of the \res{unchanged-repair-local-qwen25vl-7b-n}
pages it returned were identical to the faulty input, and \res{capped-repair-local-qwen25vl-7b-count}
of its \res{capped-repair-local-qwen25vl-7b-n} answers stopped at the 8{,}192-token output cap before
the page was complete. \Q{} changed every page it returned
(\res{unchanged-repair-local-qwen3-8b-count} of \res{unchanged-repair-local-qwen3-8b-n} unchanged),
but its edits introduced \res{regressions-local-qwen3-8b} new failures per page (CI
\res{regressions-local-qwen3-8b-ci}), most of them ARIA attributes that the element's role does not
allow. It preserved the text and structure of a share of \res{preserved-local-qwen3-8b} of the pages
(\QV: \res{preserved-local-qwen25vl-7b}; difference \res{preserved-diff-local-qwen25vl-7b-local-qwen3-8b},
CI \res{preserved-diff-local-qwen25vl-7b-local-qwen3-8b-ci}). Of \Q's pages that failed
preservation, \res{capped-repair-local-qwen3-8b-count} contained no page: the prompt and the reasoning
together filled the 32{,}768-token context window before the page was written. Most of the others
had changed the visible text.

\Q{} resolved most invalid ARIA attributes (M11,
\res{resolve-m11-local-qwen3-8b}), placeholder-only labels (M02, \res{resolve-m02-local-qwen3-8b}),
unnamed icon buttons (M04, \res{resolve-m04-local-qwen3-8b}) and removed focus indicators (M10,
\res{resolve-m10-local-qwen3-8b}), but few faults of other kinds. It almost never restored what was removed at page level, the
page title (M12, \res{resolve-m12-local-qwen3-8b}) and the page language (M06,
\res{resolve-m06-local-qwen3-8b}): none of its returned pages had a title again, and only one had a
language. Unnamed icon links (M05, \res{resolve-m05-local-qwen3-8b}) and low contrast (M03,
\res{resolve-m03-local-qwen3-8b}) also mostly remained. It scored lower on the longer W3C-derived
pages (\res{srep-w3c-local-qwen3-8b}, CI \res{srep-w3c-local-qwen3-8b-ci}) than on the authored
pages (\res{srep-authored-local-qwen3-8b}, CI \res{srep-authored-local-qwen3-8b-ci}). \QV{} resolved
only invalid ARIA attributes (M11, \res{resolve-m11-local-qwen25vl-7b}) and unnamed icon links
(M05, \res{resolve-m05-local-qwen25vl-7b}).

\textbf{Audit (C1).} \Q{} reached $S_\text{audit}={}$\res{saudit-local-qwen3-8b} (CI
\res{saudit-local-qwen3-8b-ci}), with recall \res{audit-recall-local-qwen3-8b} (CI
\res{audit-recall-local-qwen3-8b-ci}) and a precision of at least \res{audit-precision-local-qwen3-8b},
and every page gave a parseable list (parse-failure rate \res{audit-parsefail-local-qwen3-8b}).
Its hits came mostly from three operators: missing alt text (M01, \res{audit-recall-m01-local-qwen3-8b}),
placeholder-only labels (M02, \res{audit-recall-m02-local-qwen3-8b}) and \texttt{div} elements
used as buttons (M07, \res{audit-recall-m07-local-qwen3-8b}). It reported no removed page title or
page language, no unnamed icon link and no positive \texttt{tabindex} (M12, M06, M05, M09:
\res{audit-recall-m12-local-qwen3-8b} each), although the last is plain in the markup. It also
rarely reported unnamed icon buttons and removed focus indicators (M04, M10:
\res{audit-recall-m04-local-qwen3-8b} each), even though it fixed most of them when asked to
repair the same pages. What this model lists in an audit is therefore a poor guide to what it can
repair.

\QV{} found none of the 192 injected faults ($S_\text{audit}={}$\res{saudit-local-qwen25vl-7b},
recall \res{audit-recall-local-qwen25vl-7b}; difference from \Q{}
\res{saudit-diff-local-qwen25vl-7b-local-qwen3-8b}, CI
\res{saudit-diff-local-qwen25vl-7b-local-qwen3-8b-ci}), and none of the items it reported matched a
fault (precision \res{audit-precision-local-qwen25vl-7b}). Its near-greedy decoding fell into
loops that listed the same element again and again under consecutive success-criterion numbers.
On page B03-V4, for example, it reported the label of a select element three times under SC 1.4.8, 1.4.9 and
1.4.10, while the page's injected faults were a missing alt text, low contrast, a missing page
language and a demoted heading (\texttt{runs/e1/scores/audit.jsonl}).
\res{capped-audit-local-qwen25vl-7b-count} of its \res{capped-audit-local-qwen25vl-7b-n} first
answers ran into the output cap, and a share of \res{audit-parsefail-local-qwen25vl-7b} of pages
yielded no parseable list even after the JSON retry.

\textbf{Duplicate pages.} Without the two duplicate pages, $S_\text{rep}$ was
\res{srep-dedup-local-qwen3-8b} for \Q{} and \res{srep-dedup-local-qwen25vl-7b} for \QV{}, and
$S_\text{audit}$ \res{saudit-dedup-local-qwen3-8b} and \res{saudit-dedup-local-qwen25vl-7b}. No
conclusion changes.

\textbf{Answer to RQ1.} \Q{} was the stronger model in repair and audit, and the only one whose
output improved when accessibility was requested. In generation the two models could not be told
apart. Neither is ready for unattended use: few generated pages passed every check, full repairs
were rare, the better audit found fewer than a third of the known faults, and both models claimed
conformance for pages that failed checks. The profiles are also uneven (Table~\ref{tab:main}). \Q's
$S_1={}$\res{s1-local-qwen3-8b} and $S_6={}$\res{s6-local-qwen3-8b} are low, while its
$S_3={}$\res{s3-local-qwen3-8b} and $S_5={}$\res{s5-local-qwen3-8b} are high only because DCDI and
$S_3$ leave out the steps the session failed (Section~\ref{sec:rq2}).
